\documentclass{article}

\textwidth=135mm
\textheight=200mm
\oddsidemargin=20mm
\linespread{1.04}


\title{\textsc{Continuous Exponential Weights}}
\date{}
\usepackage{marginnote}
\reversemarginpar % switch default margin to the left
\usepackage{framed}
\usepackage{xcolor}
\usepackage[hypertexnames=false]{hyperref} % add backlinks to the end of bibliography items
\usepackage{nicefrac}
\usepackage{pgfplots}
\usepackage{floatrow}
\usepackage{amsmath}
\usepackage{amsthm}
\usepackage{listings}
\usepackage{tikz}
\usepackage{amssymb}
\usepackage[plain]{algorithm}
\lstset{emph={try,catch,throw,def,or,to,for,then,end,for,if,args,else,return},emphstyle=\color{blue!100!black}\textbf}
\lstset{mathescape=true,escapechar=\&,numbers=left,xleftmargin=3mm,linewidth=0.47\textwidth,framextopmargin=0pt,framexbottommargin=0pt,aboveskip=-1pt,belowskip=-1pt}
\usepackage{mdframed}
\usepackage{bm}
\usepackage{mathtools}
\usepackage[capitalise]{cleveref}
\usepackage{autonum}
\usepackage[bf]{caption}
\usepackage{graphicx}
\usepackage{natbib}
\usepackage{xspace}

\setlength{\marginparwidth}{35mm}

\definecolor{dkblue}{cmyk}{1,.54,.04,.19} 
\theoremstyle{plain}
\newtheorem{theorem}{Theorem}
\newtheorem{optprob}[theorem]{Optimisation Problem}
\newtheorem{lemma}[theorem]{Lemma}
\newtheorem{proposition}[theorem]{Proposition}
\newtheorem{corollary}[theorem]{Corollary}
\newtheorem{claim}[theorem]{Claim}
\theoremstyle{definition}
\newtheorem{definition}[theorem]{Definition}
\newtheorem{question}[theorem]{Question}
\newtheorem{fact}[theorem]{Fact}
\newtheorem{assumption}[theorem]{Assumption}
\theoremstyle{plain}
\newtheorem{exampleinner}[theorem]{Example}
\newtheorem{remark}[theorem]{Remark}

% Boxed example environment with light gray background and shadow
\mdfdefinestyle{examplestyle}{%
    backgroundcolor=gray!10,
    linewidth=0.5pt,
    linecolor=gray!50,
    shadow=true,
    shadowsize=3pt,
    shadowcolor=gray!30,
    innertopmargin=8pt,
    innerbottommargin=8pt,
    innerleftmargin=10pt,
    innerrightmargin=10pt,
    skipabove=10pt,
    skipbelow=10pt
}
\newenvironment{example}[1][]{%
    \begin{mdframed}[style=examplestyle]%
    \begin{exampleinner}[#1]%
}{%
    \end{exampleinner}%
    \end{mdframed}%
}


\newcommand{\interior}{\operatorname{int}}
\newcommand{\dom}[1]{\operatorname{dom}(#1)}
\newcommand{\ceil}[1]{\left\lceil #1 \right\rceil}
\newcommand{\R}{\mathbb R}
\newcommand{\Rd}{\mathbb{R}^d}
\newcommand{\argmin}{\operatornamewithlimits{arg\,min}}
\newcommand{\argmax}{\operatornamewithlimits{arg\,max}}
\newcommand{\ip}[1]{\langle #1 \rangle}
\newcommand{\Reg}{\mathfrak{R}}
\newcommand{\BReg}{\mathfrak{BR}}
\newcommand{\norm}[1]{\Vert #1 \Vert}
\newcommand{\E}{\mathbb E}
\newcommand{\Z}{\mathbb Z}
\newcommand{\K}{\mathcal K}
\newcommand{\cU}{\mathcal U}
\newcommand{\cN}{\mathcal N}
\newcommand{\cT}{\mathcal T}
\newcommand{\PP}{\mathbb P}
\newcommand{\Var}{\mathbb V}
\newcommand{\cE}{\mathcal E}
\newcommand{\cK}{\mathcal K}
\newcommand{\cH}{\mathcal H}
\newcommand{\cF}{\mathcal F}
\newcommand{\cO}{\mathcal O}
\newcommand{\cS}{\mathcal S}
\newcommand{\cC}{\mathcal C}
\newcommand{\cL}{\mathcal L}
\newcommand{\cG}{\mathcal G}
\newcommand{\cA}{\mathcal A}
\newcommand{\cD}{\mathcal D}
\newcommand{\cP}{\mathcal P}
\newcommand{\BS}{\mathbb S}
\newcommand{\zeros}{ \bm 0}
\newcommand{\bbP}{\mathbb P}
\newcommand{\polylog}{\operatorname{polylog}}
\newcommand{\poly}{\operatorname{poly}}
\newcommand{\const}{\operatorname{const}}
\newcommand{\conv}{\operatorname{conv}}
\newcommand{\diam}{\operatorname{diam}}
\renewcommand{\d}[1]{\operatorname{d}\!#1}
\renewcommand{\epsilon}{\varepsilon}
\newcommand{\abovelabel}[1]{\stackrel{\raisebox{0.5mm}{\textrm{\tiny \color{red} #1}}}}
\newcommand{\ta}{\tilde{a}}
\newcommand{\tpi}{\tilde{\pi}}
\newcommand{\tp}{\tilde{p}}
\newcommand{\tP}{P'}
\newcommand{\tA}{A'}
\newcommand{\fmd}{\texttt{MD}}
\newcommand{\tfmd}{\widetilde{\texttt{MD}}}

\newcommand{\mdalg}{\textsf{MD}\xspace}
\newcommand{\cewaalg}{\textsf{CEWA}\xspace}
\newcommand{\smdalg}{\textsf{SCMD}\xspace}
\newcommand{\scewaalg}{\textsf{SCEWA}\xspace}
\newcommand{\expbyopt}{\textsf{ExpByOpt}\xspace}
\newcommand{\olo}{\texttt{OLO}\xspace}
\newcommand{\blo}{\texttt{BLO}\xspace}
\newcommand{\bco}{\texttt{BCO}\xspace}
\newcommand{\oco}{\texttt{OCO}\xspace}
\newcommand{\hf}{\hat{\ell}\xspace}
    
\begin{document}
\maketitle

In this note, we aim to improve the dimension dependence of bandit convex optimization (\bco).
The main theoretical tool we use is the analysis of a continuous version of mirror descent (\mdalg).
Then we move on to \bco.
The current state of affairs can be summarized as follows:
\begin{enumerate}
    \item \oco\footnote{Online Convex Optimization} reduces to \olo\footnote{Online Linear Optimization}: $f(x) - f(x_\star) \leq \langle f'(x), x - x_\star \rangle$.
    \item \textbf{Random} algs are \textit{suboptimal} for \texttt{OCO} and \textit{pointless} for \texttt{OLO}.
    \item First order information, $f'(x)$, makes \texttt{BCO} as hard as \texttt{OLO}.
    \item \textbf{Random} algs are \textit{necessary} for \texttt{BCO} and \texttt{BLO}.
    \item \textbf{Random} algs are \textit{free} for \texttt{BLO}\footnote{Bandit Linear Optimization}
    \item \textbf{Random} algs appear to be \textit{expensive} for \texttt{BCO}: upper bound is $d^{2.5}\sqrt{n}$ while lower bound is $d\sqrt{n}$.
\end{enumerate}
Our main goal here is to prove a $d^{1.5} \sqrt{n}$ regret upper bound for \bco,
which partially closes the gap between the upper and lower bounds for \bco.

\newpage
\tableofcontents
\newpage

\section{Continuous exponential weights}
We study a repeated game between a player and an adversary, which proceeds in $n$ rounds.
Let $\K \subset \R^d$ be a compact convex subset with non-empty interior, and the set $\cF \subseteq \mathsf{C}(\K)$ be the class of continuous functions defined on $\K$.
In each round $t \in [n]$, the player chooses an \textbf{action} $A_t \in \K$ (possibly randomly) and incurs a loss of $\ell_t(A_t)$, where $\ell_t \in \cF$ is a loss function chosen by the adversary, simultaneously with the player.
Throughout this section, we study the \textit{full information} setting where the adversary reveals the whole loss function $\ell_t$ to the player at the end of each round.

The player is allowed to randomize its actions, by choosing a distribution $\pi_t$ over $\K$ and sampling $A_t \sim \pi_t$.
The distribution $\pi_t$ is called the \textbf{policy}.
Precisely, in round $t$, the player chooses a distribution $\pi_t \in \mathsf{P}(\K)$ where
$\mathsf{P}(\K)$ is the set of all probability measures supported on $\K$.
The goal of the player is to minimize regret with respect to a comparator action $a_\star \in \K$, defined as
\begin{align}
    \Reg_n(a_\star) \triangleq \sum_{t=1}^{n} \ell_t(A_t) - \ell_t(a_\star)\,.
\end{align}
Note that $\Reg_n(a_\star)$ is a random variable,
and the goal of the player is to minimize this quantity, either in expectation or with high probability.

\subsection{The algorithm and its regret bound}
The Continuous Exponential Weights Algorithm (\cewaalg), shown in \cref{alg:MD},
is a natural extension of the discrete exponential weights algorithm to the continuous setting.
\begin{algorithm}[H]
    \begin{minipage}{12cm}
        \begin{mdframed}
            \begin{lstlisting}
args: $\eta > 0,  \pi_1$
for $t = 1$ to $n$:
    Play $A_t \sim \pi_t$ and observe $\ell_t$
    For $x \in \K$, let $\tpi_{t+1}(x) = \pi_t(x) \exp\left(-\eta \ell_t(x)\right)$
    Set $\pi_{t+1}(x) = \frac{\tpi_{t+1}(x)}{\int_\K \tpi_{t+1}(y) d\mu(y)}$
\end{lstlisting}
            \caption{
                \cewaalg algorithm.
            }\label{alg:MD}
        \end{mdframed}
    \end{minipage}
\end{algorithm}
\noindent
Before stating the regret bound for \cewaalg, we need to introduce some notations.
Define the set of probability densities supported on $\K$ as
\begin{align}
    \cP(\K)
    \triangleq
    \left\{
    p:\K\to\R_+
    \;\middle|\;
    \textit{measurable}, \int_\K p(x)\d\mu(x)=1
    \right\}.
\end{align}
From now on, we will be mainly concerned with policies that are absolutely continuous with respect to the Lebesgue measure $\mu$ on $\K$, so that we can identify a policy $\pi$ with its density function,
for which we abuse the notation and also write $\pi$ to denote the density function of the policy $\pi$.
The \cewaalg algorithm is a special case of the mirror descent algorithm with \textbf{the negative entropy} as the potential function, which is defined as
\begin{align}
    H(v) \triangleq \int_{\K} v(x) \log v(x) - v(x) + 1 d\mu(x)\,, \qquad \text{with } 0\log 0 \triangleq 0\,,
\end{align}
for a measurable function $v:\K \to \R_+$.
The Bregman divergence associated with the negative entropy is the generalized Kullback-Leibler divergence, defined as
\begin{align}
    D_H(p,q)
    =
    \int_\K p(x)\log\frac{p(x)}{q(x)} - p(x) + q(x)\d x\,,
\end{align}
for $p,q$ being two measurable functions from $\K$ to $\R_+$.
For a probability density $p \in \cP(\K)$, and a loss function $\ell \in \cF$, define the dual pairing as
\begin{align}
    \langle \ell, p \rangle
    =
    \int_\K \ell(x) p(x) d\mu(x)\,,
\end{align}
which is always finite since $\ell$ is continuous on a compact set and $p$ is a probability density.
Finally, for a comparator fixed policy $\pi_\star \in \cP(\K)$, define the \textbf{pseudo-regret} with respect to $\pi_\star$ as
\begin{align}
    \Reg_n(\pi_\star) \triangleq
    \sum_{t=1}^{n} \langle \ell_t, \pi_t - \pi_\star \rangle\,.
\end{align}

The following theorem gives a regret bound for the regret of \cewaalg algorithm.
Note that the comparator in the following regret bound is a density $\pi_\star$ instead of a point $a_\star$.
However, these two types of comparators are essentially equivalent since the class of losses we study are continuous functions on a compact set, and the regret with respect to a Dirac measure $\delta_{a_\star}$ can be approximated arbitrarily well by the regret with respect to a density $\pi_\star$ that is supported on a small neighborhood of $a_\star$.
\begin{theorem}\label{thm:MD}
    Let $\pi_1 \in \cP(\K)$ be an arbitrary initial policy with $\pi_1>0$ everywhere on $\K$,
    and $\pi_\star \in \cP(\K)$ be an arbitrary comparator policy,
    such that $\pi_1 \ll \pi_\star$ (otherwise the upper bound is trivial).
    Then the regret of \cewaalg with respect to $\pi_\star$ satisfies
    \begin{align}
        \Reg_n(\pi_\star)
         & \leq
        \frac{1}{\eta}
        \left(
        D_H(\pi_\star, \pi_1)
        +
        \sum_{t=1}^{n} D_H(\pi_t, \tilde{\pi}_{t+1})
        \right)\,.
    \end{align}
\end{theorem}

\begin{proof}
    We have
    \begin{align}
        \Reg_n(\pi_\star)
         & =
        \sum_{t=1}^{n} \langle \ell_t, \pi_t - \pi_\star \rangle \\
         & =
        \sum_{t=1}^{n}
        \underbrace{\langle \ell_t, \pi_t - \pi_{t+1} \rangle}_{\triangleq \texttt{st}_t \text{ or stability}}
        +
        \underbrace{\langle \ell_t, \pi_{t+1} - \pi_\star \rangle}_{\triangleq \texttt{tr}_t \text{ or tracking}}\,.
    \end{align}
    It is standard that the normalized exponential update for $\pi_{t+1}$ can be equivalently written as the following optimization problem
    \begin{align}
        \pi_{t+1}
         & =
        \argmin_{\pi \in \cP(\K)}
        \eta \langle \ell_t, \pi \rangle + D_H(\pi,\pi_t)\,,
    \end{align}
    then the first order optimality condition for this optimization problem gives us
    \begin{align}
         & \forall \bar{\pi} \in \cP(\K), \quad
        \left\langle
        \eta \ell_t + \log \pi_{t+1} - \log \pi_t,
        \bar{\pi} - \pi_{t+1}
        \right\rangle
        \geq 0                                  \\
        \Rightarrow
         &
        \texttt{tr}_t
        =
        \langle \ell_t, \pi_{t+1} - \pi_\star \rangle
        \leq
        \frac{1}{\eta}
        \left\langle
        \log \pi_{t+1} - \log \pi_t,
        \pi_\star - \pi_{t+1}
        \right\rangle\,.
    \end{align}
    The RHS in the last display can be written in terms of the Bregman divergences as
    \begin{align}
        \left\langle
        \log \pi_{t+1} - \log \pi_t,
        \pi_\star - \pi_{t+1}
        \right\rangle
         & =
        D_H(\pi_\star, \pi_t)
        -
        D_H(\pi_{t+1}, \pi_t)
        -
        D_H(\pi_\star, \pi_{t+1})\,,
    \end{align}
    which gives us a bound on the tracking terms
    \begin{align}
        \sum_{t=1}^{n} \texttt{tr}_t
         & =
        \sum_{t=1}^{n}
        \langle \ell_t, \pi_{t+1} - \pi_\star \rangle \\
         & \leq
        \frac{1}{\eta}
        \left(
        \sum_{t=1}^{n}
        D_H(\pi_\star, \pi_t)
        -
        D_H(\pi_\star, \pi_{t+1})
        -
        D_H(\pi_{t+1}, \pi_t)
        \right)                                       \\
         & =
        \frac{1}{\eta}
        \left(
        D_H(\pi_\star, \pi_1)
        -
        \underbrace{D_H(\pi_\star, \pi_{n+1})}_{\geq 0}
        -
        \sum_{t=1}^{n} D_H(\pi_{t+1}, \pi_t)
        \right)
        \\
         & \leq
        \frac{1}{\eta}
        \left(
        D_H(\pi_\star, \pi_1)
        -
        \sum_{t=1}^{n} D_H(\pi_{t+1}, \pi_t)
        \right)
        \,.
    \end{align}
    To analyze the stability term $\texttt{st}_t$,
    note that $\textsf{CEWA}$ stores the latest loss in the unnormalized density $\tilde{\pi}_{t+1}$, as
    \begin{align}
        \log \tilde{\pi}_{t+1}
        =
        \log \pi_t - \eta \ell_t,
        \qquad\text{so}\qquad
        \ell_t
        =
        \frac{1}{\eta}
        \left(
        \log \pi_t - \log \tilde{\pi}_{t+1}
        \right).
    \end{align}
    Therefore, the stability term can be written as
    \begin{align}
        \texttt{st}_t
         & =
        \langle \ell_t, \pi_t - \pi_{t+1} \rangle \\
         & =
        \frac{1}{\eta}
        \left\langle
        \log \pi_t - \log \tilde{\pi}_{t+1},
        \pi_t - \pi_{t+1}
        \right\rangle                             \\
         & =
        \frac{1}{\eta}
        \left(
        D_H(\pi_{t+1}, \pi_t)
        +
        D_H(\pi_t, \tilde{\pi}_{t+1})
        -
        D_H(\pi_{t+1}, \tilde{\pi}_{t+1})
        \right).
    \end{align}
    Thus, by dropping the negative term $-D_H(\pi_{t+1}, \tilde{\pi}_{t+1})$,
    we have proved the second inequality,
    \begin{align}
        \Reg_n(\pi_\star)
         & \leq
        \frac{1}{\eta}
        \left(
        D_H(\pi_\star, \pi_1)
        +
        \sum_{t=1}^{n} D_H(\pi_t, \tilde{\pi}_{t+1})
        \right).
    \end{align}
\end{proof}

\subsection{A discussion on generalization to unbounded losses}
Since, $\cF \subseteq \mathsf{C}(\K)$ and continuous functions on compact sets are automatically bounded, the bilinear form $\langle f, p \rangle$ is finite for every $p \in \cP(\K)$.
However, one might wonder about the \textbf{largest class of measurable loss functions} that are compatible with continuous exponential weights algorithm.
Firstly, the next theorem characterizes the class of functions $v$ for which $H(v)$ is finite, i.e. $\dom{H}$.
\begin{theorem}
    \label{thm:negentropy-LlogLspace}
    Define
    \begin{align}
        \Phi(x) \triangleq \begin{cases}
                               (1+x) \log(1+x) - x & x \ge 0\,, \\
                               +\infty             & x < 0\,.
                           \end{cases}
    \end{align}
    Then, the domain of $H$ is given by
    \begin{align}
        \dom{H} = L\log\!L(\K) \triangleq \left\{v: \K \to \R_+ \middle| \text{measurable}, \int_\K \Phi(v(x))d\mu(x) < +\infty\right\}\,. \label{eq:LlogLspace_definition}
    \end{align}
\end{theorem}
The $L\log\!L(\K)$ in \cref{eq:LlogLspace_definition} is the Orlicz space associated with the function $\Phi$.
The dual space of $L\log\!L(\K)$ is the Orlicz space associated with the convex conjugate $\Phi^\star$ of $\Phi$, which is given by
\begin{align}
    \Phi^\star(x) = \begin{cases}
                        e^x - x - 1, & x \ge 0  \\
                        0,           & x < 0\,,
                    \end{cases}
\end{align}
and
\begin{align}
    L\!\exp(\K) \triangleq \left\{f: \K \to \R \middle| \textit{measurable}, \int_\K \Phi^\star(f(x)) d\mu(x) < +\infty\right\}\,.
\end{align}
Then, the dual pairing (bilinear form) between functions $p \in L\!\log\!L(\K)$ and $f \in L\!\exp(\K)$ is finite and given by
\begin{align}
    \langle f, p \rangle \triangleq \int_\K f(x) p(x) d\mu(x) < +\infty\,.
\end{align}
The curious fact here is that the set $L\!\exp(\K)$ coincides with the set of measurable functions $f$ such that the denominator $\int_\K \exp(-\eta f(x)) d\mu(x)$ in the update rule of \cewaalg is finite for every $\eta > 0$.
I still don't have a good intuition for why this is the case.



\newpage
\section{Stochastic continuous exponential weights}
In this section, we study the bandit version of the online optimization problem.
The main difference is that the player only observes the loss incurred by its action,
namely $\ell_t(A_t)$, instead of the entire loss function $\ell_t$.
\textbf{Unlike the full information setting, the sequence $\ell_t$ must be fixed before the game starts.}
Essentially, \textit{the regret is not really well-defined if $\ell_t$ is chosen adaptively by the adversary after observing the previous actions of the player}\footnote{The adversary can decide to show only $\ell_t(\cdot)=1$ for $t\geq2$, if $a_1\neq a_\star$, and otherwise show $\ell_t(\cdot)=0$ for all $t$ when $a_1 = a_\star$, turning the whole game a trivial needle in a haystack problem.}.

The main difficulty is that the player cannot form $\tpi_t$, since it does not observe the entire loss function $\ell_t$.
The key idea is to use the sampled action $A_t$ together with the observed loss $\ell_t(A_t)$ to construct an \textit{estimated loss function} $\hf_t$.
This is the first modification to the \cewaalg algorithm, leading to the Stochastic Exponential Weights Algorithm (\scewaalg) in \cref{alg:stochastic-cewa}.
A second modification is that the algorithm need not play according to $\pi_t$ itself; instead it samples $A_t \sim \hat{\pi}_t$, where $\hat{\pi}_t$ is a distribution constructed from $\pi_t$.

Philosophically, the resulting algorithm has very little in common with the mirror descent algorithm.
Precisely, $\pi_t$ is not really playing the role of the `policy' in round $t$, but rather it plays the role of a `sufficient statistic' that summarizes the history of the game up to time $t$, and the player uses it to construct the \textit{playing policy} $\hat{\pi}_t$ and the \textit{estimator map} $\cE_t: \K \times \R \to \mathsf{C}(\K)$ that maps the observed action and loss to an estimated loss function $\hf_t$.

\begin{algorithm}[h!]
    \begin{minipage}{13cm}
        \begin{mdframed}
            \begin{lstlisting}
args: $\eta > 0$, $\pi_1$
for $t = 1$ to $n$:
    Choose $\hat{\pi}_t$ and $\cE_t$ based on $\pi_t$
    Play $A_t \sim \hat{\pi}_t$ and observe $Y_t = \ell_t(A_t)$
    Set $\hf_t = \cE_t(A_t, Y_t)$
    For $x \in \K$, let $\tpi_{t+1}(x) = \pi_t(x) \exp\left(-\eta \hf_t(x)\right)$
    Set $\pi_{t+1}(x) = \frac{\tpi_{t+1}(y)}{\int_\K \tpi_{t+1}(y) d\mu(y)}$
\end{lstlisting}
            \caption{\scewaalg algorithm.}\label{alg:stochastic-cewa}
        \end{mdframed}
    \end{minipage}
\end{algorithm}

\begin{proposition}
    \label{prop:cexp_regret_bound}
    Let $\pi_1 \in \cP(\K)$ be an arbitrary initial policy,
    and $\pi_\star \in \cP(\K)$ be an arbitrary comparator policy.
    The expected-regret of \scewaalg with respect to $\pi_\star$ satisfies
    \begin{align}
        \E[
            \Reg_T(\pi_\star)
        ]
         & \leq
        \frac{1}{\eta}D_H(\pi_\star,\pi_1)
        +
        \sum_{t=1}^T
        \E\left[
            \langle \ell_t,\hat{\pi}_t-\pi_\star\rangle
            -
            \langle \hf_t,\pi_t-\pi_\star\rangle
            +
            \frac{1}{\eta}
            D_{H}
            \left(
            \pi_t,
            \tilde{\pi}_{t+1}
            \right)
            \right].
    \end{align}
\end{proposition}

\begin{proof}
    Starting from the definition of regret, and using that conditionally on the past
    $A_t\sim \hat{\pi}_t$, we can write
    \begin{align}
        \E[\Reg_T(\pi_\star)]
         & =
        \E\left[
            \sum_{t=1}^T
            \ell_t(A_t) - \ell_t(A_\star)
            \right]
         & (\text{where } A_\star \sim \pi_\star)
        \\
         & =
        \E\left[
            \sum_{t=1}^T
            \langle \ell_t,\hat{\pi}_t-\pi_\star\rangle
        \right]                                   \\
         & =
        \E\left[
            \sum_{t=1}^T
            \left(
            \langle \ell_t,\hat{\pi}_t-\pi_\star\rangle
            -
            \langle \hf_t,\pi_t-\pi_\star\rangle
            \right)
            +
            \sum_{t=1}^T
            \langle \hf_t,\pi_t-\pi_\star\rangle
            \right].
    \end{align}
    Applying \cref{thm:MD} with $\ell_t$ replaced by $\hf_t$, we get
    \begin{align}
        \sum_{t=1}^T
        \langle \hf_t,\pi_t-\pi_\star\rangle
         & \leq
        \frac{1}{\eta}D_H(\pi_\star,\pi_1)
        +
        \frac{1}{\eta}
        \sum_{t=1}^T
        D_H(\pi_t,\tilde{\pi}_{t+1}).
    \end{align}
    Plugging this into the previous display gives the result.
\end{proof}

By Fenchel duality for the negative entropy,
\begin{align}
    D_H(\pi_t,\tilde{\pi}_{t+1})
     & =
    D_{H^\star}
    \left(
    H'(\tilde{\pi}_{t+1}),
    H'(\pi_t)
    \right) =
    D_{H^\star}
    \left(
    H'(\pi_t)-\eta \hf_t,
    H'(\pi_t)
    \right),
\end{align}
so another way to write the regret bound is
\begin{align}
    \E[\Reg_T(\pi_\star)]
     & \leq
    \frac{1}{\eta}D_H(\pi_\star,\pi_1)
    +
    \sum_{t=1}^T
    \E\left[
        \langle \ell_t,\hat{\pi}_t-\pi_\star\rangle
        -
        \langle \hf_t,\pi_t-\pi_\star\rangle
        +
        \frac{1}{\eta}
        D_{H^\star}
        \left(
        H'(\pi_t)-\eta \hf_t,
        H'(\pi_t)
        \right)
        \right]\,,\label{eq:regret_bound_cexp_with_Bregman_divergence}
\end{align}
which has the benefit of isolating only the terms concerning time step $t$.

\subsection{The one step complexity \texorpdfstring{$\Lambda$}{Lambda}}
We still haven't answered the question of how to choose $\hat{\pi}_t$ and $\cE_t$ in \cref{alg:stochastic-cewa}.
An answer to this question is given by the regret bound in \cref{prop:cexp_regret_bound}, which suggests that one should choose $\hat{\pi}_t$ and $\cE_t$ to minimize the one-step quantity
\begin{align}
    \E\left[
        \langle \ell_t,\hat{\pi}_t-\pi_\star\rangle
        -
        \langle \hf_t,\pi_t-\pi_\star\rangle
        +
        \frac{1}{\eta}
        D_{H^\star}
        \left(
        H'(\pi_t)-\eta \hf_t,
        H'(\pi_t)
        \right)
        \right].
\end{align}
Note that when the learner is choosing the playing policy $\hat{\pi}_t$ and the estimator $\cE_t$, the learning rate $\eta$ and the current \cewaalg policy $\pi_t$ are given to the learner, while the loss function $\ell_t$ and the comparator policy $\pi_\star$ are unknown to the learner.
Motivated by this observation, define the one step complexity as
\begin{align}
    \Lambda_{\pi,\eta}(\hat{\pi},\cE \| \pi_\star, \ell)
    \triangleq
    \frac{1}{\eta}
    \E_{X\sim\hat{\pi}}\left[
        \langle \ell,\delta_X-\pi_\star\rangle
        -
        \langle \cE_X^\ell,\pi-\pi_\star\rangle
        +
        \frac{1}{\eta}
        D_{H^\star}
        \left(
        H'(\pi)-\eta \cE_X^\ell,
        H'(\pi)
        \right)
        \right],
\end{align}
where the randomness is with respect to $X \sim \hat{\pi}$ and $\cE_X^\ell = \cE(X,\ell(X))$.
The reason for the extra $\frac{1}{\eta}$ factor in the definition of $\Lambda$ will become clear shortly.

Suppose that there exists a choice of $\hat{\pi}_t$ and $\cE_t$ such that $\Lambda_{\pi_t,\eta}(\hat{\pi}_t,\cE_t \| \pi_\star, \ell_t) \leq \Lambda_\star$ for every $t$, \textbf{and any $\eta>0$}, then the regret bound in \cref{prop:cexp_regret_bound} can be bounded as
\begin{align}
    \E[\Reg_T(\pi_\star)]
     & \leq
    \frac{1}{\eta}D_H(\pi_\star,\pi_1)
    +
    \eta
    \sum_{t=1}^T
    \Lambda_{\pi_t,\eta}(\hat{\pi}_t,\cE_t \| \pi_\star, \ell_t) \\
     & \leq
    \frac{1}{\eta}D_H(\pi_\star,\pi_1)
    +
    \eta T \Lambda_\star                                         \\
     & \leq
    2\sqrt{D_H(\pi_\star,\pi_1) \Lambda_\star T}\,,
\end{align}
where the last inequality is setting $\eta = \sqrt{\frac{D_H(\pi_\star,\pi_1)}{T\Lambda_\star}}$.
Note that the extra $\frac{1}{\eta}$ factor in the definition of $\Lambda$ gives the correct scaling of $\Lambda$ for $O(\sqrt{T})$ regret bound.


It can be the case that one can only ensure $\Lambda_{\pi_t,\eta}(\hat{\pi}_t,\cE_t \| \pi_\star, \ell_t) \leq \Lambda_\star L(\eta)$ where $L(\eta) = \polylog(1/\eta)$.
In that case, the regret bound becomes
\begin{align}
    \E[\Reg_T(\pi_\star)]
     & \leq
    \frac{1}{\eta}D_H(\pi_\star,\pi_1)
    +
    \eta T \Lambda_\star L(\eta) \\
     & \leq
    2\sqrt{D_H(\pi_\star,\pi_1) \Lambda_\star T} \sqrt{L\left(\sqrt{\frac{D_H(\pi_\star,\pi_1)}{T\Lambda_\star}}\right)}\,,
\end{align}
which is still $O(\sqrt{T})$ up to polylogarithmic factors, for essentially the same choice of $\eta$ as before.
Further refinement in terms of the choice of $\eta$ might be possible depending on the exact form of $L(\eta)$.


\newpage
\section{The convex losses and the optimistic kernel}
In this section, we construct a specific choice of $\hat{\pi}_t$ and $\cE_t$,
for the case where \textbf{the loss functions $\ell_t$ are convex and bounded}.
Let
\begin{align}
    \cF_{\mathrm{cvx}} \triangleq \left\{f:\K\to[0, 1] \middle| \,\, f \text{ is convex}\right\}\,.
\end{align}
Let $\ell \in \cF_{\mathrm{cvx}}$, and divide the one step complexity $\Lambda_{\pi,\eta}(\hat{\pi},\cE \| \pi_\star, \ell)$ into a bias term and a variance term:
\begin{align}
    \Lambda_{\pi,\eta}(\hat{\pi},\cE \| \pi_\star,\ell)
     & =
    \underbrace{
        \frac{1}{\eta}
        \E_{X\sim\hat{\pi}}\left[
            \langle \ell,\hat{\pi}-\pi_\star\rangle
            -
            \langle \cE_X^\ell,\pi-\pi_\star\rangle
            \right]
    }_{\triangleq \Lambda^{\mathsf{b}}_{\eta,\pi}(\cE,\hat{\pi}\| \pi_\star,\ell)}
    +
    \underbrace{
        \frac{1}{\eta^2}
        \E_{X\sim\hat{\pi}}\left[
            D_{H^\star}
            \left(
            H'(\pi)-\eta \cE_X^\ell,
            H'(\pi)
            \right)
            \right]
    }_{\triangleq \Lambda^{\mathsf{v}}_{\eta,\pi}(\cE,\hat{\pi}\|\ell)}\,.
\end{align}

First, we restrict ourselves to the case where $\cE_X^\ell$ is \textit{multiplicative} in the observed loss $\ell(X)$, namely for every $x\in\K$ and $\ell\in\cF_{\mathrm{cvx}}$, let
\begin{align}
    \cE_x^\ell = \cE(x,\ell(x)) = \ell(x) s_x - \ell(x) \mathbf{1}\,,
\end{align}
for some function $s_x:\K\to\cF$, where $\cF\subseteq\mathsf{C}(\K)$ is the class of \textit{estimated loss functions}.
The reason for the extra $-\ell(x) \mathbf{1}$ term will be clear later in the analysis of the variance term.
Note that $\cE_x^\ell$ is not necessarily a convex function.

The goal of the learner with the choice of $\cE$ in round $t$ is to estimate the whole loss function $\ell_t$ from the single point observation $(A_t, \ell_t(A_t))$, in expectation over the randomness of $A_t$.
This task is difficult in general, but such estimators are known for the case of $\ell$ being a linear function.
In the next section, we explain how such an estimator can be constructed for the case of convex losses, by using the kernel-based construction method of \cite{bubeck2016kernelbasedmethodsbanditconvex}.
It is important to note that the resulting estimator is not unbiased, but it is an \textit{optimistic} estimator, in the sense that it is a lower bound on the true loss function $\ell$.

To ease the notation, we refer to $\hat{\pi} = p$ and $\pi = q$ in the rest of the section.

\subsection{The optimistic kernel}
Let $\pi \in \cP(\K)$ be a given policy, and $\lambda \in (0,1)$ be a free parameter.
Define $\tilde{\ell}$ as
\marginpar{We are hiding the dependence of $\tilde{\ell}$ on $\lambda$ and $p$ to ease the notation.}
\begin{align}
    \tilde{\ell}(t) \triangleq \E_{X\sim p}\left[
        \left(1 - \frac{1}{\lambda}\right)\ell(X) + \frac{1}{\lambda} \ell\left((1-\lambda) X + \lambda t\right)
        \right]\,.
\end{align}
First, we show that $\left(1 - \frac{1}{\lambda}\right)\ell(X) + \frac{1}{\lambda} \ell\left((1-\lambda) X + \lambda t\right)$ is a lower bound on $\ell(t)$ for every $t\in\K$.
To see this, note that by convexity of $\ell$, we have
\begin{align}
    \ell((1-\lambda) X + \lambda t)
     & \leq
    (1-\lambda) \ell(X) + \lambda \ell(t)\,,
\end{align}
which gives us
\begin{align}
    \left( 1 - \frac{1}{\lambda}\right)\ell(X) + \frac{1}{\lambda} \ell\left((1-\lambda) X + \lambda t\right)
     & \leq
    \ell(t)\,,
\end{align}
which proves that $\tilde{\ell}(t) \leq \ell(t)$ for every $t\in\K$.
Also note that $\tilde{\ell}$ is a convex function since it is a non-negative linear combination of convex functions;
however, it is not necessarily bounded in $[0,1]$ even if $\ell$ is bounded in $[0,1]$.

Next, we show that, for any $p \in \cP(\K)$, there exists an unbiased one-point estimator of $\tilde{\ell}$.
Precisely, we show that there exists a function $s_x$ such that $\tilde{\ell}(t) = \E_{X\sim p}[\ell(X) s_X(t)]$ for every $t\in\K$.
From the definition of $\tilde{\ell}$, we have
\begin{align}
    \tilde{\ell}(t)
     & = \int_{\Rd}
    p(x) \left[
        \left(1 - \frac{1}{\lambda}\right) \ell(x)
        + \frac{1}{\lambda} \ell\left(
        (1-\lambda) x + \lambda t
        \right)
    \right]\,dx     \\
     & =
    \left(1 - \frac{1}{\lambda}\right)
    \int_{\Rd}
    p(x)
    \ell(x)
    dx
    +
    \frac{1}{\lambda}
    \int_{\Rd}
    p(x)
    \ell\left(
    (1-\lambda) x + \lambda t
    \right) \,dx    \\
     & =
    \left(1 - \frac{1}{\lambda}\right)
    \int_{\Rd}
    p(x)
    \ell(x)
    dx
    +
    \frac{1}{\lambda (1-\lambda)^d}
    \int_{\Rd}
    p\left(
    \frac{y - \lambda t}{1 - \lambda}
    \right)
    \ell(y) \,dy    \\
     & =
    \int_{\Rd}
    p(x)
    \ell(x)
    \left(
    1 - \frac{1}{\lambda}
    +
    \frac{1}{\lambda (1-\lambda)^d}
    \cdot
    p\left(
    \frac{x - \lambda t}{1 - \lambda}
    \right)
    p(x)^{-1}
    \right) \,dx\,,
\end{align}
where we have used the change of variable $y = (1-\lambda)x + \lambda t$ in the third equality.
Therefore, let
\begin{align}
    s_x(t) \triangleq \left(1 - \frac{1}{\lambda} + \frac{1}{\lambda (1-\lambda)^d} \cdot p\left(\frac{x - \lambda t}{1 - \lambda}\right) p(x)^{-1}\right)\,,
\end{align}
which gives us
\begin{align}
    \tilde{\ell}(t) =  \E_{X\sim p}[\ell(X) s_X(t)]
    = \E_{X\sim p}[\cE_X^\ell(t) + \ell(X)]
    = \E_{X\sim p}[\cE_X^\ell(t)] - \langle \ell, p \rangle
    \,,
\end{align}
for every $t\in\K$.

By plugging in this choice into the definition of the bias term, we have
\begin{align}
    \Lambda^{\mathsf{b}}_{\eta,\pi}(\cE,\hat{\pi}\| \pi_\star,\ell)
     & =
    \frac{1}{\eta}
    \E_{X\sim p}\left[
        \langle \ell,\delta_X -\pi_\star\rangle
        -
        \langle \cE_X^\ell,q-\pi_\star\rangle
    \right] \\
     & =
    \frac{1}{\eta}
    \E_{X\sim p}\left[
        \langle \ell,\delta_X -\pi_\star\rangle
        -
        \langle \ell(X) s_X - \ell(X)\cdot \mathbf{1},q-\pi_\star\rangle
    \right] \\
     & =
    \frac{1}{\eta}
    \E_{X\sim p}\left[
        \langle \ell,\delta_X -\pi_\star\rangle
        -
        \langle \ell(X) s_X ,q-\pi_\star\rangle
    \right] \\
     & =
    \frac{1}{\eta}
    \left(
    \langle \ell,p\rangle
    -
    \langle \tilde{\ell},q\rangle
    +
    \langle \tilde{\ell} - \ell, \pi_\star\rangle
    \right) \\
     & \leq
    \frac{1}{\eta}
    \left(
    \langle \ell,p\rangle
    -
    \langle \tilde{\ell},q\rangle
    \right)\,,
\end{align}
where the inequality follows from the fact that $\tilde{\ell} \leq \ell$ pointwise.

Next, we show a choice of $p$ that makes the upper bound on the bias term vanish,
and thus the bias term itself will be non-positive.
Precisely, we show that if $p$ is the law of the random vector
\begin{align}
    X \triangleq \lambda \sum_{k=0}^\infty (1-\lambda)^k T_k\,,
\end{align}
where $T_k$ are i.i.d. samples from $q$, then
\begin{align}
    \langle \ell,p\rangle
     & =
    \langle \tilde{\ell},q\rangle\,.
    \label{eq:core_vanishes_bias}
\end{align}
The law of the random vector $X$, which we denote by $p$, is also called the \textit{core} of the distribution $q$ with parameter $\lambda$.
To see why \cref{eq:core_vanishes_bias} holds, write
\begin{align}
    \langle \tilde{\ell},q\rangle
     & =
    \int_{\Rd}
    \tilde{\ell}(t) q(t) \d t \\
     & =
    \int_{\Rd}
    \left(
    \int_{\Rd}
    p(x) \left[
        \left(1 - \frac{1}{\lambda}\right)\ell(x) + \frac{1}{\lambda} \ell\left((1-\lambda) x + \lambda t\right)
        \right] \d x
    \right)
    q(t) \d t                 \\
     & =
    \left(1 - \frac{1}{\lambda}\right)
    \int_{\Rd}
    p(x)
    \ell(x) \d x
    + \frac{1}{\lambda}
    \int_{\Rd}
    \int_{\Rd}
    p(x) q(t)
    \ell\left((1-\lambda) x + \lambda t\right)
    \d t \d x                 \\
     & =
    \left(1 - \frac{1}{\lambda}\right)
    \langle \ell,p\rangle
    +
    \frac{1}{\lambda}
    \E_{X\sim p,\,T\sim q}
    \left[
        \ell\left((1-\lambda)X+\lambda T\right)
    \right]                   \\
     & =
    \left(1 - \frac{1}{\lambda}\right)
    \langle \ell,p\rangle
    +
    \frac{1}{\lambda}
    \E
    \left[
        \ell
        \left(
        \lambda T+\lambda\sum_{k=0}^{\infty}(1-\lambda)^{k+1}T_k
        \right)
    \right]                   \\
     & =
    \left(1 - \frac{1}{\lambda}\right)
    \langle \ell,p\rangle
    +
    \frac{1}{\lambda}
    \E
    \left[
        \ell
        \left(
        \lambda\sum_{k=0}^{\infty}(1-\lambda)^k T'_k
        \right)
    \right]                   \\
     & =
    \left(1 - \frac{1}{\lambda}\right)
    \langle \ell,p\rangle
    +
    \frac{1}{\lambda}
    \langle \ell,p\rangle     \\
     & =
    \langle \ell,p\rangle\,,
\end{align}
where we used the fact that if $T\sim q$ is independent of
$X\sim p$, then $(1-\lambda)X+\lambda T\sim p$ by the distributional identity
\begin{align}
    (1-\lambda)X+\lambda T
    \stackrel{d}{=}
    \lambda\sum_{k=0}^{\infty}(1-\lambda)^k T'_k\,,
\end{align}
where $(T'_k)_{k\geq0}$ are again i.i.d. samples from $q$.

Therefore, with this choice of $p$, the bias term $\Lambda^{\mathsf{b}}_{\eta,q}(\cE,p\| \pi_\star,\ell)$ vanishes, and we are left with only the variance term as
\begin{align}
    \Lambda_{q,\eta}(p,\cE \| \pi_\star,\ell)
     & =
    \Lambda^{\mathsf{b}}_{\eta,q}(\cE,p\| \pi_\star,\ell)
    +
    \Lambda^{\mathsf{v}}_{\eta,q}(\cE,p\|\ell)
    \\
     & \leq
    \Lambda^{\mathsf{v}}_{\eta,q}(\cE,p\|\ell)\,,
\end{align}
which is independent of the comparator policy $\pi_\star$.



\newpage
\section{Bounding the variance}
Throughout this section we fix the estimator $\cE$ and the policy $p$ as defined in the previous section, and we want to upper bound the variance term $\Lambda^{\mathsf{v}}_{\eta,q}(\cE,p\|\ell)$ uniformly over $\ell \in \cF_{\mathrm{cvx}}$.
Recall the \scewaalg variance term defined as
\begin{align}
    \Lambda^{\mathsf{v}}_{\eta, q}(\cE, p \| \ell)
     & =
    \frac{1}{\eta^2}
    \E_{X\sim p}\left[
        D_{H^\star}\left(
        H'(q) - \eta \cE^\ell_X, H'(q)
        \right)
        \right]
    \,.
\end{align}
In terms of $\eta$ dependency, we need a $\polylog(\frac{1}{\eta})$ upper bound on this quantity to get a $\tilde{O}(\sqrt{T})$ regret bound.


First, we write the variance in a simpler form.
For $\eta \ge 0, u \in \R$, define
\begin{align}
    \phi_\eta(u) \triangleq e^{-\eta u} + \eta u - 1\,,
\end{align}
which is approximately quadratic $\eta^2 u^2/2$ when $\eta u$ is small.
\begin{lemma}
    \label{prop:Lambda_in_terms_of_phi}
    Let $X \sim p$ and $T \sim q$, then
    \begin{align}
        \Lambda^{\mathsf{v}}_{\eta, q}(\cE, p \| \ell) & =
        \frac{1}{\eta^2}
        \E\left[
            \phi_\eta(\cE_X^\ell(T))
            \right]\,.
        \label{eq:Lambda_in_terms_of_phi}
    \end{align}
\end{lemma}
The proof is deferred to \cref{proof:Lambda_in_terms_of_phi}.

The next lemma shows that for this choice of estimator (or basically any multiplicative estimator), the variance term is maximized when $\ell(x) = 1$ for all $x$, which is a consequence of strong convexity of the function $x \mapsto e^{-x} + x - 1$ and the fact that it is minimized at $x=0$.
\begin{proposition}[The worst case is $\ell \equiv 1$]
    \label{lem:appendix_f_one_reduction}
    We have
    \begin{align}
        \sup_{\ell: \K \to [0,1]}
        \Lambda^{\mathsf{v}}_{\eta, q}(\cE, p \| \ell)
        \leq
        \Lambda^{\mathsf{v}}_{\eta, q}(\cE, p \| 1)\,.
    \end{align}
\end{proposition}
The proof is deferred to \cref{proof:appendix_f_one_reduction}, and note that \textbf{even convexity of the loss functions is not needed for this reduction}.
This is very convenient because using \cref{prop:Lambda_in_terms_of_phi}, we have
\begin{align}
    \sup_{\ell: \K \to [0,1]}
    \Lambda^{\mathsf{v}}_{\eta, q}(\cE, p_X\| \ell)
     & =
    \frac{1}{\eta^2}
    \E\left[
        \phi_\eta(\cE_X^1(T))
        \right]
    =
    \frac{1}{\eta^2}
    \E\left[
        \phi_\eta(s_X(T) - 1)
        \right]
    \,. \label{eq:variance_reduction_to_f_one}
\end{align}

\subsection{The relative density $R$ and the $f$-divergence}
Next, we show that $s_x(t)$ can be written in terms of the relative density $r(x,t)$ defined as
\marginpar{We use multiple notations to refer to same densities: $p_X=p$ and $p_T=p_{T_0}=q$.}
\begin{align}
    r(x, t) & \triangleq  \frac{p_{X, T_0}(x, t)}{p_{T_0}(t) p_X(x)}\,, \quad\text{then}\quad
    s_x(t) = 1 + \frac{r(x, t) - 1}{\lambda}\,.
    \label{eq:s_in_terms_of_r}
\end{align}
To this end, write
\begin{align}
    X
    =
    \lambda T_0 + (1-\lambda)X',
    \qquad
    X' \stackrel{d}{=} X,
    \qquad
    X' \perp T_0 .
\end{align}
\marginpar{We implicitly extend all densities to $\Rd$ by zero extension.}
Therefore, conditionally on $T_0=t$, the density of $X$ is
\begin{align}
    p_{X \mid T_0=t}(x)
    =
    \frac{1}{(1-\lambda)^d}
    p\left(
    \frac{x-\lambda t}{1-\lambda}
    \right)\,.
\end{align}
Thus, starting from the definition of $s_x(t)$, we can write
\begin{align}
    s_x(t)
     & =
    1 - \frac{1}{\lambda} + \frac{1}{\lambda (1-\lambda)^d} \cdot p\left(\frac{x - \lambda t}{1 - \lambda}\right) p(x)^{-1}
    \\
     & =
    1 - \frac{1}{\lambda} + \frac{1}{\lambda} \frac{p_{X|T_0=t}(x)}{p(x)}
    \\
     & =
    1 - \frac{1}{\lambda} + \frac{1}{\lambda} \frac{p_{X, T_0}(x, t)}{p_X(x) p_{T_0}(t)}
    \\
     & =
    1-\frac{1}{\lambda}
    +
    \frac{1}{\lambda} r(x,t)
    =
    1+\frac{r(x,t)-1}{\lambda}\,,
\end{align}
which proves \cref{eq:s_in_terms_of_r}.

Motivated by this we define two related random variables $R$ and $R_0$ as
\begin{align}
    R   \triangleq r(X, T) = \frac{p_{X, T_0}(X, T)}{p_X(X) p_{T_0}(T)}\,,
    \qquad
    R_0 \triangleq r(X, T_0) = \frac{p_{X, T_0}(X, T_0)}{p_X(X) p_{T_0}(T_0)}\,,
\end{align}
where $T \sim q$ is a fresh independent sample from $q$.
Observe that $R$ and $R_0$ satisfy the size biasing relation $\E[g(R_0)] = \E[g(R) R]$ for every measurable function $g$.

Let
\begin{align}
    \alpha = \frac{\eta}{\lambda}\,, \qquad \text{and}\qquad f_\alpha(u) = \phi_\alpha(u-1)\,,
\end{align}
and note that $f_\alpha$ is a standard $f$-divergence generator, i.e. it is convex and satisfies $f_\alpha(1)=f_\alpha'(1)=0$.
The next proposition shows that the variance term can be written in terms of the $f$-divergence generated by $f_\alpha$.
\begin{proposition}
    \label{prop:f_eta_lambda_properties}
    We have
    \begin{align}
        \sup_{\ell: \K \to [0,1]}
        \Lambda^{\mathsf{v}}_{\eta, q}(\cE, p_X\| \ell)
         & =\frac{1}{\eta^2}
        \E\left[ f_{\alpha}(R) \right]
        \,.
        \label{eq:phi_eta_f_divergence_relation}
    \end{align}
\end{proposition}
The proof is deferred to \cref{proof:f_eta_lambda_properties}.
Therefore, everything boils down to
\begin{align}
    \sup_{\ell: \K \to [0,1]}
    \Lambda^{\mathsf{v}}_{\eta, q}(\cE, p_X\| \ell)
     & =
    \frac{1}{\eta^2} D_{f_{\alpha}}(\PP_{X, T_0} || \PP_X \otimes \PP_{T_0})\,.
\end{align}
The question, then, is how close $\PP_{X, T_0}$ and $\PP_X \otimes \PP_{T_0}$ are in terms of $f_{\alpha}$-divergence.
We need them to become closer at a rate of $\alpha^2 \polylog(\frac{1}{\alpha})$.

We can pause here to write the expected regret of the \scewaalg algorithm with the optimistic kernel and
policy $p$, using \cref{eq:regret_bound_cexp_with_Bregman_divergence}, as
\begin{align}
    \E[\Reg_T(\pi_\star)]
     & \leq
    \frac{1}{\eta}D_H(\pi_\star,\pi_1)
    +
    \E\left[
        \sum_{t=1}^T
        \E\left[
            \langle \ell_t,\hat{\pi}_t-\pi_\star\rangle
            -
            \langle \hf_t,\pi_t-\pi_\star\rangle
            +
            \frac{1}{\eta}
            D_{H^\star}
            \left(
            H'(\pi_t)-\eta \hf_t,
            H'(\pi_t)
            \right)
            \middle|
            \pi_t
            \right]
        \right]
    \\
     & =
    \frac{1}{\eta}D_H(\pi_\star,\pi_1)
    +
    \E\left[
        \sum_{t=1}^T
        \eta
        \Lambda_{\pi_t,\eta}(\hat{\pi}_t,\cE_t \| \pi_\star, \ell_t)
        \right]
    \\
     & \leq
    \frac{1}{\eta}D_H(\pi_\star,\pi_1)
    +
    \E\left[
        \sum_{t=1}^T
        \eta
        \Lambda^{\mathsf{v}}_{\pi_t, \eta}(\hat{\pi}_t, \cE_t \| \ell_t)
        \right]
    \label{eq:regret_bound_cexp_variance_only}
    \\
     & \leq
    \frac{1}{\eta}D_H(\pi_\star,\pi_1)
    +
    \E\left[
        \sum_{t=1}^T
        \eta
        \Lambda^{\mathsf{v}}_{\pi_t, \eta}(\hat{\pi}_t, \cE_t \| \mathbf{1})
        \right]
    \\
     & =
    \frac{1}{\eta}D_H(\pi_\star,\pi_1)
    +
    \frac{1}{\eta}
    \E\left[
    \sum_{t=1}^T
    D_{f_{\alpha}}(\PP_{A_t, T_0^{(t)}} || \PP_{A_t} \otimes \PP_{T_0^{(t)}})
    \right]
    \\
     & =
    \frac{1}{\eta}D_H(\pi_\star,\pi_1)
    +
    \frac{1}{\eta}
    \E\left[
    \sum_{t=1}^T
    I_{f_{\alpha}}(A_t; T_0^{(t)})
    \right]\,.
\end{align}



\newpage
\section{Bounding the $f_\alpha$ divergence}
In this section, we explore ways to upper bound $D_{f_{\alpha}}(\PP_{X, T_0} || \PP_X \otimes \PP_{T_0})$.
The objective is to find the right set of assumptions on $q=\PP_{T_0}$ (that can be ensured throughout the algorithmic trajectory whp) that lead to the desired $\alpha^2 \polylog(\frac{1}{\alpha})$ bound.

\subsection{The KL route}
The next proposition shows that $\PP_{X, T_0}$ and $\PP_X \otimes \PP_{T_0}$ are actually very close in terms of KL divergence.
\begin{proposition}
    \label{prop:appendix_exact_kl_mi}
    It holds that
    \begin{align}
        D_{\mathrm{KL}}(\PP_{X,T_0}\,\|\,\PP_X \otimes \PP_{T_0})
        = d\log\left(\frac{1}{1-\lambda}\right).
        \label{eq:appendix_exact_kl_mi}
    \end{align}
\end{proposition}
The proof is deferred to \cref{proof:appendix_exact_kl_mi}.
Note that the bound in \eqref{eq:appendix_exact_kl_mi} is essentially a constant when $\lambda=\tilde{\Theta}(1/d)$,
which says that these two distributions are very similar, at least in terms of KL divergence.

\noindent
While the KL bound is a constant, \textbf{it doesn't give the $\alpha^2$ scaling} that we need,
but it gives a blueprint
\begin{align}
    D_{f_{\eta, \lambda}}(\PP_{X, T_0} || \PP_X \otimes \PP_{T_0})
     & \stackrel{(?)}{\lesssim}
    \alpha^2 \polylog(\frac{1}{\alpha}) \cdot \poly(\frac{1}{\lambda}) \cdot D_{\mathrm{KL}}(\PP_{X,T_0}\,\|\,\PP_X \otimes \PP_{T_0})\,.
    \label{eq:blueprint}
\end{align}
To make the comparison more meaningful, we multiply the KL generator by a factor of $\alpha^2$ which gives
\begin{align}
    k_\alpha(u) = \alpha^2 \left(u\log(u) - (u - 1)\right)\,,
\end{align}
which matches the value of $f_\alpha$ and its first and second derivatives at $u=1$:
\begin{align}
    f_\alpha(1)   & = k_\alpha(1) = 0\,,          \\
    f_\alpha'(1)  & = k_\alpha'(1) = 0\,,         \\
    f_\alpha''(1) & = k_\alpha''(1) = \alpha^2\,.
\end{align}
We get the desired bound if we can show that
\begin{align}
    \E\left[
        f_\alpha(R)
        \right]
    \stackrel{(?)}{\lesssim}
    \polylog(\frac{1}{\alpha}) \cdot \poly(\frac{1}{\lambda}) \cdot
    \E\left[k_\alpha(R)\right]\,.
\end{align}
This inequality is not true in general, so answering this question requires a better understanding of the distribution of $R$.

\subsection{The chi-squared route}
Every $f$-divergence in its natural form (after affine transformation) is positive, convex and satisfies $g(1) = 0$, so it is upper bounded by its second order Taylor expansion around $1$.
Basically,
\begin{align}
    g(u) & \leq \frac{g''(1)}{2} (u - 1)^2 + O((u - 1)^3)\,,
\end{align}
where $(u-1)^2$ is the generator of the $\chi^2$-divergence.
As long as one can show that $R$ is concentrated around $1$, then the higher order terms are negligible, and the $f$-divergence can be upper bounded by a constant times the $\chi^2$-divergence, which is in turn upper bounded by a constant times $\E[R^2] - 1$.

Motivated by this, we show that a bound on $\E[R^2]$ suffices.
\begin{proposition}
    \label{prop:appendix_f_chi_squared_bound}
    For every $\alpha > 0$ and every $u \geq 0$, we have
    \begin{align}
        D_{f_{\alpha}}(\PP_{X, T_0}\|\PP_X \otimes \PP_{T_0}) & \leq
        \frac{1}{2}
        e^{\alpha}
        \alpha^2
        D_{\chi^2}(\PP_{X, T_0}\|\PP_X \otimes \PP_{T_0})\,.
    \end{align}
\end{proposition}
\begin{proof}
    Observe that by Taylor's theorem with Lagrange remainder, for every $u \geq 0$, we have
    \begin{align}
        f_\alpha(u) = \phi_\alpha(u-1) & \leq \frac{\alpha^2}{2} (u - 1)^2 e^{-\alpha(u'-1)}\,, \qquad\text{for some } u' \in [0, u]\,.
    \end{align}
    Since $\alpha > 0$, we have $e^{-\alpha(u'-1)} \leq e^{\alpha}$, so
    \begin{align}
        \phi_\alpha(u-1) & \leq \frac{\alpha^2}{2} e^{\alpha} (u - 1)^2\,.
        \label{eq:f_eta_lambda_chi_squared_bound}
    \end{align}
    Using this we get a bound in terms of $\chi^2$-divergence as
    \begin{align}
        f_{\alpha}(R) & =
        \phi_\alpha(R - 1)
        \leq
        \frac{1}{2}
        \alpha^2
        e^{\alpha}
        (R - 1)^2\,,
    \end{align}
    integrating both sides gives the desired bound.
\end{proof}
How loose is the bound?
The larger $u'$ is, the looser the bound is, and $u'$ can be as large as $R$, which can be arbitrarily large.

The next result shows that the exact quadratic scaling
of $D_{f_\alpha}(P\|Q)$ in $\alpha$ as $\alpha \downarrow 0$ holds if and only if $D_{\chi^2}(P\|Q)<\infty$.
\begin{proposition}
    \label{prop:exact_quadratic_forces_chi2}
    For every two densities $P$ and $Q$ such that $P \ll Q$, we have
    \begin{align}
        \liminf_{\alpha\downarrow0}
        \frac{1}{\alpha^2}
        D_{f_{\alpha}}(P\|Q)
        =
        \frac{1}{2}
        \chi^2(P\|Q)\,.
    \end{align}
\end{proposition}
\begin{proof}
    For every fixed $u\ge0$,
    \begin{align}
        \lim_{\alpha\downarrow0}
        \frac{\phi_\alpha(u-1)}{\alpha^2}
        =
        \lim_{\alpha\downarrow0}
        \frac{e^{-\alpha(u-1)} - 1 + \alpha(u-1)}{\alpha^2}
        =
        \frac12(u-1)^2\,.
    \end{align}
    Hence,
    \begin{align}
        \liminf_{\alpha \downarrow 0}
        \frac{D_{f_{\alpha}}(P\|Q)}{\alpha^2}
        =
        \liminf_{\alpha\downarrow0}
        \E\left[
            \frac{\phi_{\alpha}(R-1)}{\alpha^2}
            \right]
        \geq
        \E\left[
            \liminf_{\alpha\downarrow0}
            \frac{\phi_{\alpha}(R-1)}{\alpha^2}
            \right]
        =
        \E\left[\frac{(R-1)^2}{2}\right]\,,
    \end{align}
    where the inequality follows from Fatou's lemma.
    For the upper bound, from \cref{eq:f_eta_lambda_chi_squared_bound} we have
    \begin{align}
        \phi_\alpha(u-1)
        \le
        \frac12 e^{\alpha} \alpha^2(u-1)^2,
        \qquad u\ge0.
    \end{align}
    Therefore, we have
    \begin{align}
        \lim_{\alpha\downarrow0}
        \frac{1}{\alpha^2}
        D_{f_\alpha}(P\|Q)
         & =
        \lim_{\alpha\downarrow0}
        \frac{1}{\alpha^2}
        \E\left[
            \phi_\alpha(R-1)
            \right]
        \le
        \lim_{\alpha\downarrow0}
        \frac{1}{2 e^{\alpha}}
        \E[(R-1)^2]
        =
        \frac12 \E[(R-1)^2]\,.
    \end{align}
    and rearranging and taking limit gives the result.
\end{proof}

Perhaps the most interesting fact about chi-squared reduction is that it gives the desired result
if one just assumes that all iterates $\pi_t$ are Gaussian:
\begin{theorem}[Gaussian kernels have explicit $\chi^2$ control]
    \label{thm:gaussian_chi_squared_control}
    Assume that $q=\cN(m,\Sigma)$ on $\Rd$, where $\Sigma \succ 0$, and let
    $T_0,T_1,\ldots$ be i.i.d. samples from $q$. Define
    \begin{align}
        X
        =
        \lambda\sum_{k=0}^{\infty}(1-\lambda)^kT_k,
        \qquad
        \lambda\in(0,1).
    \end{align}
    Then $X$ is Gaussian and
    \begin{align}
        X
        \sim
        \cN\left(
        m,
        \frac{\lambda}{2-\lambda}\Sigma
        \right).
    \end{align}
    Moreover,
    \begin{align}
        D_{\chi^2}
        \left(
        \PP_{X,T_0}
        \middle\|
        \PP_X\otimes\PP_{T_0}
        \right)
        =
        \frac{1}{(1-\lambda)^{2d}}-1.
        \label{eq:gaussian_chi_squared_explicit}
    \end{align}
\end{theorem}
\noindent
The proof is somewhat classical and is dropped for now.

In particular, taking $\alpha=\eta/\lambda$ and $\lambda=1/d$ gives
\begin{align}
    \sup_{\ell:\K\to[0,1]}
    \Lambda^{\mathsf{v}}_{\eta,q}(\cE,p_X\|\ell)
     & \le
    \frac{e^{\eta/\lambda}}{2\lambda^2}
    \left(
    \frac{1}{(1-\lambda)^{2d}}-1
    \right) \\
     & =
    \frac12 e^{\eta d} d^2
    \left(
    \frac{1}{(1-\frac{1}{d})^{2d}}-1
    \right) \\
     & \leq
    \frac{2e-1}{2} e^{\eta d} d^2\,,
\end{align}
which by putting into \cref{eq:regret_bound_cexp_variance_only} gives
\begin{align}
    \E[\Reg_T(\pi_\star)]
     & \leq
    \frac{1}{\eta}D_H(\pi_\star,\pi_1)
    +
    \E\left[
        \sum_{t=1}^T
        \eta
        \Lambda^{\mathsf{v}}_{\pi_t, \eta}(\hat{\pi}_t, \cE_t \| \ell_t)
        \right]
    \\
     & \leq
    \frac{1}{\eta}D_H(\pi_\star,\pi_1)
    +
    \frac{2e-1}{2} T\eta e^{\eta d} d^2\,,
\end{align}
which by choosing $\eta = \sqrt{\frac{2D_H(\pi_\star,\pi_1)}{(2e-1)Td^2}}$ gives
\begin{align}
    \E[\Reg_T(\pi_\star)]
     & \leq
    \sqrt{2(2e-1)D_H(\pi_\star,\pi_1)Td^2} \cdot e^{\sqrt{\frac{2D_H(\pi_\star,\pi_1)}{(2e-1)T}}}\,,
\end{align}


\clearpage
\bibliographystyle{plainnat}
\bibliography{../all}
\clearpage

\appendix
\section{Appendix}
\label{sec:appendix}

\subsection{Proofs for $\Lambda_{\mathrm{var}}$ bounds}
\subsubsection{Proof of \cref{prop:Lambda_in_terms_of_phi}}
\label{proof:Lambda_in_terms_of_phi}
\begin{proof}
    Note that $H^\star(v) = \int_\K e^{v(t)} dt$, and $H'(q) = \log(q)$.
    Then, starting from the definition of the variance term, we have
    \begin{align}
        \Lambda^{\mathsf{v}}_{\eta, q}(\cE, p \| \ell)
         & =
        \frac{1}{\eta^2}
        \E\left[
            D_{H^\star}\left(
            H'(q) - \eta \cE^\ell_X, H'(q)
            \right)
            \right]
        \\
         & =
        \frac{1}{\eta^2}
        \E\left[
            H^\star(H'(q) - \eta \cE^\ell_X) - H^\star(H'(q)) + \eta \langle \cE^\ell_X, q \rangle
            \right]
        \\
         & =
        \E\left[
            \int_\K \left(e^{-\eta \cE_X^\ell(t)} - 1 + \eta \cE_X^\ell(t)\right) q(t) dt
            \right]
        \\
         & =
        \E\left[
            \phi_\eta(\cE_X^\ell(T))
            \right]
        \,,
    \end{align}
    where the third equality follows from the fact that
    \begin{align}
        H^\star(H'(q)) & = \int_\K e^{\log(q(t))} dt
        = \int_\K q(t) dt = 1\,.\qedhere
    \end{align}
\end{proof}

\subsubsection{Proof of \cref{lem:appendix_f_one_reduction}}
\label{proof:appendix_f_one_reduction}
\begin{proof}
    Note that from \cref{prop:Lambda_in_terms_of_phi}, we have
    \begin{align}
        \Lambda^{\mathsf{v}}_{\eta, q}(\cE, p \| \ell) & =
        \frac{1}{\eta^2}
        \E\left[
            \phi_\eta(\cE_X^\ell(T))
            \right]
        =
        \frac{1}{\eta^2}
        \E\left[
            \phi_\eta(\ell(X)(s_X(T) - 1))
            \right]\,,
    \end{align}
    where the second equality follows from the definition of $\cE_X^\ell(T) = \ell(X)s_X(T) - \ell(X)$.
    We claim that for every $a \in [0,1]$ and every $u \in \R$,
    \begin{align}
        \phi_\eta(au) \leq a \phi_\eta(u) \leq \phi_\eta(u).
        \label{eq:appendix_phi_contraction}
    \end{align}
    Once this is proved, we can apply it with $a=\ell(x)$ and $u=(s_x^\lambda(t)-1)$.
    Let
    \begin{align}
        g(x) = e^{-x} + x - 1\,,
    \end{align}
    which is convex, minimized at $x=0$ and satisfies $g(0)=0$, so for every $x \in \R$ and every $a \in [0,1]$ we have
    \begin{align}
        g(ax) \leq a g(x) + (1-a)g(0) = a g(x).
    \end{align}
    Finally, just observe that $\phi_\eta(u) = g(\eta u)$ gives us the first inequality in \cref{eq:appendix_phi_contraction}, and the second inequality is trivial.
\end{proof}

\subsubsection{Proof of \cref{prop:f_eta_lambda_properties}}
\label{proof:f_eta_lambda_properties}
\begin{proof}
    Starting from \cref{eq:variance_reduction_to_f_one}, we have
    \begin{align}
        \sup_{\ell: \K \to [0,1]}
        \Lambda^{\mathsf{v}}_{\eta, q}(\cE, p \| \ell)
        =
        \Lambda^{\mathsf{v}}_{\eta, q}(\cE, p \| \mathbf{1})
         & =
        \frac{1}{\eta^2}
        \E\left[
            \phi_\eta(s_X(T) - 1)
            \right]
        =
        \frac{1}{\eta^2}
        \E\left[
            \phi_\eta\left(
            \lambda^{-1}(R - 1)
            \right)
            \right]\,,
        \label{eq:appendix_f_eta_lambda_reduction}
    \end{align}
    where the second equality follows from \cref{eq:s_in_terms_of_r} and the definition of $R$.
    Now, by the definition of $\phi_\eta(u) = e^{-\eta u} + \eta u - 1$,
    \begin{align}
        \phi_\eta\left(
        \lambda^{-1}(R - 1)
        \right)
         & =
        e^{-\frac{\eta}{\lambda}(R - 1)}
        + \frac{\eta}{\lambda} (R - 1) - 1
        = f_{\alpha}(R)\,,
    \end{align}
    where $\alpha = \eta/\lambda$ and the last equality follows from the definition of $f_\alpha$.
\end{proof}

\subsubsection{Proof of \cref{prop:appendix_exact_kl_mi}}
\label{proof:appendix_exact_kl_mi}
\begin{proof}
    The identity
    \begin{align}
        D_{\mathrm{KL}}(\PP_{X,T_0}\,\|\,\PP_X \otimes \PP_{T_0}) = I(T_0;X)
    \end{align}
    is the standard definition of mutual information.
    So it remains only to compute $I(T_0;X)$.
    Using the decomposition
    \begin{align}
        X = \lambda T_0 + (1-\lambda)X',
    \end{align}
    with $X'$ independent of $T_0$ and $X' \stackrel{d}{=} X$, we have
    \begin{align}
        I(T_0;X)
         & = h(X) - h(X\mid T_0).
        \label{eq:appendix_mi_entropy_start}
    \end{align}
    Now condition on $T_0=t$. Then
    \begin{align}
        X\mid(T_0=t) = \lambda t + (1-\lambda)X'.
    \end{align}
    The term $\lambda t$ is a deterministic translation, and differential entropy is invariant under translation, hence
    \begin{align}
        h(X\mid T_0=t) = h\left((1-\lambda)X'\right).
    \end{align}
    Next, differential entropy scales according to
    \begin{align}
        h(aY) = h(Y) + d\log|a|,
    \end{align}
    for every nonzero scalar $a$.
    Applying this with $a=1-\lambda$ gives
    \begin{align}
        h\left((1-\lambda)X'\right)
        = h(X') + d\log(1-\lambda).
    \end{align}
    Since $X' \stackrel{d}{=} X$, we have $h(X') = h(X)$. Therefore,
    \begin{align}
        h(X\mid T_0=t)
        = h(X) + d\log(1-\lambda),
    \end{align}
    and this does not depend on $t$. Averaging over $T_0$ yields
    \begin{align}
        h(X\mid T_0)
        = h(X) + d\log(1-\lambda).
    \end{align}
    Substituting this into \eqref{eq:appendix_mi_entropy_start},
    \begin{align}
        I(T_0;X)
         & = h(X) - \left(h(X) + d\log(1-\lambda)\right)
        = d\log\left(\frac{1}{1-\lambda}\right),
    \end{align}
    which proves \eqref{eq:appendix_exact_kl_mi}.
\end{proof}


\end{document}
